\section{The finite-full-shift semiring skeleton}
\label{sec:semiring}

Let $A_n$ be an $n$-element alphabet and let
\[
F_n=A_n^{\mathbb Z}
\]
denote the corresponding two-sided full shift, considered up to topological
conjugacy.  The source objects are two-sided because their entropy and
product structure are convenient and canonical.  The verifier built later
is a separate one-sided countable Markov shift.

\begin{definition}[Alphabet product and alphabet-sum]
For $m,n\ge1$, set
\[
F_m\boxtimes F_n:=F_{A_m\times A_n}\cong F_{mn},
\qquad
F_m\boxplus F_n:=F_{A_m\sqcup A_n}\cong F_{m+n}.
\]
We call $\boxplus$ the \emph{alphabet-sum}, or alphabet coproduct followed by
the full-shift functor.
\end{definition}

This terminology is deliberate.  The dynamical disjoint union
$F_m\sqcup F_n$ permits a point to stay in one component forever.  By
contrast, $F_{A_m\sqcup A_n}$ permits arbitrary temporal mixing between the
two alphabets.  We therefore make no claim that $\boxplus$ is a categorical
coproduct in a category of subshifts.

Alphabet cardinality immediately gives a semiring skeleton:
\[
[F_m]\boxtimes[F_n]=[F_{mn}],\qquad
[F_m]\boxplus[F_n]=[F_{m+n}].
\]
Both operations are invariant under alphabet relabeling.  The intrinsic norm
used below is topological entropy,
\[
h(F_n)=\log n,
\qquad
h(F_m\boxtimes F_n)=h(F_m)+h(F_n).
\]

We also freeze an additive order and successor:
\[
F_a<F_b
\quad\Longleftrightarrow\quad
F_a\boxplus F_c=F_b\ \text{for some }c\ge1,
\qquad
S(F_d)=F_d\boxplus F_1=F_{d+1}.
\]
Thus comparison, multiplication, and increment can all be stated in source
language.  In an executable presentation, the witness $F_c$ is exposed by
successor rather than supplied by an existential edge guard.

\begin{remark}[Structural credit and its limit]
The construction does not discover a new semiring law: it chooses a simple
family of symbolic objects whose alphabet cardinalities reproduce
$\N_{\ge1}$.  This is nevertheless source-intrinsic enough to earn the
structural arithmetic coordinate used in \cref{sec:route}.  The later Euler
identity is credited separately at the determinant coordinate.
\end{remark}

The key question is what happens when these instructions are made temporal.
The next section turns them into a local deterministic graph and carefully
avoids the one forbidden shortcut: an edge that already knows whether a
factor exists.
